\chapter{Conclusion et Perspectives}

\section{Synthèse du projet}

Le projet \textit{HR Analytics - Employee Attrition \& Performance} a permis de concevoir une solution de Business Intelligence complète, capable de transformer des données RH brutes en un outil d'aide à la décision stratégique. En suivant une méthodologie rigoureuse, nous avons franchi les étapes de l'analyse des besoins, de la modélisation dimensionnelle en étoile, du processus ETL et enfin de la visualisation interactive.

L'analyse a révélé que l'attrition n'est pas un phénomène aléatoire mais qu'elle est étroitement liée à des facteurs identifiables tels que la charge de travail (heures supplémentaires), l'équilibre vie-travail et la structure de rémunération. La mise en place de cet entrepôt de données offre désormais à l'organisation une vision claire et objective de son capital humain, permettant de passer d'une gestion intuitive à une gestion basée sur les faits.

\section{Limites et Recommandations}

Bien que robuste, le système actuel présente certaines limites inhérentes à la nature statique des données utilisées. Pour maximiser l'impact des résultats obtenus, nous recommandons à la direction des ressources humaines :
\begin{itemize}
    \item \textbf{Ciblage des politiques de rétention :} Mettre en place des programmes de bien-être spécifiques pour les départements à forte pression (Ventes) et revoir la politique des heures supplémentaires.
    \item \textbf{Flexibilité :} Encourager le télétravail pour les collaborateurs éloignés géographiquement afin de réduire l'attrition liée aux déplacements.
\end{itemize}

\section{Perspectives d'évolution}

Ce projet constitue une première étape solide vers une gestion mature des ressources humaines par les données. Plusieurs perspectives d'évolution peuvent être envisagées pour enrichir cette solution :

D'une part, l'intégration de nouvelles sources de données, telles que les données issues des réseaux sociaux d'entreprise ou les feedbacks qualitatifs des entretiens de sortie, permettrait d'affiner la compréhension des motifs psychologiques de départ.

D'autre part, la transition du \textit{Descriptive Analytics} vers le \textit{Predictive Analytics} représente l'évolution naturelle de ce système. L'implémentation de modèles de Machine Learning (tels que les forêts aléatoires ou les réseaux de neurones) permettrait de calculer un "score de risque de départ" individuel. Cette approche prédictive offrirait aux managers la possibilité d'intervenir en amont, par le biais d'actions de rétention personnalisées, avant même que la décision de départ ne soit formalisée par l'employé.
